\ifdefined\gkver\else\def\gkver{student}\fi
\documentclass[\gkver]{gaokaozhenti}
\begin{document}

\chapter{2018年11月浙江}

\begin{enumerate}
\item
%% number: 1
%% paperName: 2018年11月浙江选考第 1 题 3 分
%% typeId: 1
%% type: 单选题
%% chapter: 电路
%% point: 电流强度
%% method: 无
%% score: 3
%% degree: 950
%% duplicateId: 0
%% body:
下列物理量属于标量的是\xzanswer{C}
\fourchoices[answer=C]
{速度}
{加速度}
{电流}
{电场强度}
%% answer: C
%% memo:
\memoanswer{
加速度、速度、电场强度既有大小又有方向，为矢量；虽然电流有方向，但只有一个方向，没有正负之分，所以为标量，C 正确。
}

\item
%% number: 2
%% paperName: 2018年11月浙江选考第 2 题 3 分
%% typeId: 1
%% type: 单选题
%% chapter: 其他
%% point: 物理学史
%% method: 无
%% score: 3
%% degree: 950
%% duplicateId: 0
%% body:
发现电流磁效应的物理学家是\xzanswer{B}
\fourchoices[answer=B]
{法拉第}
{奥斯特}
{库仑}
{安培}
%% answer: B
%% memo:
\memoanswer{
奥斯特发现了电流的磁效应，法拉第发现了电磁感应现象，库仑发现了库仑定律，安培发现了分子电流假说，B 正确。
}

\item
%% number: 3
%% paperName: 2018年11月浙江选考第 3 题 3 分
%% typeId: 1
%% type: 单选题
%% chapter: 运动和力
%% point: 力学单位制
%% method: 无
%% score: 3
%% degree: 650
%% duplicateId: 0
%% body:
用国际单位制的基本单位表示电场强度的单位，下列正确的是\xzanswer{D}
\fourchoices[answer=D]
{$\UNC$}
{$\UVm$}
{$\Ukg \cdot \Um/(\UC \cdot \Us^{2})$}
{$\Ukg \cdot \Um/(\UA \cdot \Us^{3})$}
%% answer: D
%% memo:
\memoanswer{
电场强度 $E = \frac{F}{q}$，电场力的单位为 $\UN$，电荷量的单位为 $\UC$，所以电场强度的单位是 $\UNC$；而
$1\UNC = 1\frac{\Ukg \cdot \Umsq}{\UA \cdot \Us} = 1\Ukg \cdot \Um/(\UA \cdot \Us^{3})$，D 正确。
}

\item
%% number: 4
%% paperName: 2018年11月浙江选考第 4 题 3 分
%% typeId: 1
%% type: 单选题
%% chapter: 直线运动
%% point: 运动图象
%% method: 无
%% score: 3
%% degree: 850
%% duplicateId: 0
%% body:
一辆汽车沿平直道路行驶，其 $v-t$ 图象如图所示。在 $t=0$ 到 $t=40\Us$ 这段时间内，汽车的位移是\xzanswer{C}

\onepicture[w=6.0cm]{figs/04.png}

\fourchoices[answer=C]
{0}
{$30\Um$}
{$750\Um$}
{$1200\Um$}
%% answer: C
%% memo:
\memoanswer{
在 $v-t$ 图象中，图线与时间轴围成的面积表示位移，故在 $40\Us$ 内的位移为
$x = \frac{1}{2} \times (10 + 40) \times 30\Um = 750\Um$，C 正确。
}

\item
%% number: 5
%% paperName: 2018年11月浙江选考第 5 题 3 分
%% typeId: 1
%% type: 单选题
%% chapter: 功和能
%% point: 动能势能
%% method: 无
%% score: 3
%% degree: 850
%% duplicateId: 0
%% body:
奥运会比赛项目撑杆跳高如图所示，下列说法不正确的是\xzanswer{B}

\onepicture[w=4.5cm]{figs/05.png}

\fourchoices[answer=B]
{加速助跑过程中，运动员的动能增加}
{起跳上升过程中，杆的弹性势能一直增加}
{起跳上升过程中，运动员的重力势能增加}
{越过横杆后下落过程中，运动员的重力势能减少动能增加}
%% answer: B
%% memo:
\memoanswer{
加速助跑过程中速度增大，动能增加，A 正确；撑杆从开始形变到撑杆恢复形变时，先是运动员的部分动能转化为杆的弹性势能，后弹性势能转化为运动员的动能与重力势能，杆的弹性势能不是一直增加，B 错误；起跳上升过程中，运动员的高度不断增大，重力势能增加，C 正确；越过横杆下落过程中，高度降低、速度增大，重力势能转化为动能，即重力势能减少、动能增加，D 正确。
}

\item
%% number: 6
%% paperName: 2018年11月浙江选考第 6 题 3 分
%% typeId: 1
%% type: 单选题
%% chapter: 电场
%% point: 电场线
%% method: 无
%% score: 3
%% degree: 750
%% duplicateId: 0
%% body:
等量异种电荷的电场线如图所示，下列表述正确的是\xzanswer{C}

\onepicture[w=5.5cm]{figs/06.png}

\fourchoices[answer=C]
{a 点的电势低于 b 点的电势}
{a 点的场强大于 b 点的场强，方向相同}
{将一负电荷从 a 点移到 b 点电场力做负功}
{负电荷在 a 点的电势能大于在 b 点的电势能}
%% answer: C
%% memo:
\memoanswer{
沿电场线方向电势降低，故 a 点电势高于 b 点电势，A 错误；电场线的疏密程度表示电场强度大小，电场线越密电场强度越大，故 a 点的场强大于 b 点的场强；电场线的切线方向为场强方向，故 a、b 两点的电场强度方向不同，B 错误；负电荷在低电势处电势能大，将负电荷从 a 点（高电势）移到 b 点（低电势），电势能增大，电场力做负功，C 正确、D 错误。
}

\item
%% number: 7
%% paperName: 2018年11月浙江选考第 7 题 3 分
%% typeId: 1
%% type: 单选题
%% chapter: 磁场
%% point: 左手定则
%% method: 无
%% score: 3
%% degree: 550
%% duplicateId: 0
%% body:
电流天平是一种测量磁场力的装置，如图所示。两相距很近的通电平行线圈 Ⅰ 和 Ⅱ，线圈 Ⅰ 固定，线圈 Ⅱ 置于天平托盘上。当两线圈均无电流通过时，天平示数恰好为零。下列说法正确的是\xzanswer{A}

\onepicture[w=8.0cm]{figs/07.png}

\fourchoices[answer=A]
{当天平示数为负时，两线圈电流方向相同}
{当天平示数为正时，两线圈电流方向相同}
{线圈 Ⅰ 对线圈 Ⅱ 的作用力大于线圈 Ⅱ 对线圈 Ⅰ 的作用力}
{线圈 Ⅰ 对线圈 Ⅱ 的作用力与托盘对线圈 Ⅱ 的作用力是一对相互作用力}
%% answer: A
%% memo:
\memoanswer{
当两线圈电流方向相同时相互吸引，方向相反时相互排斥，故当天平示数为正时，两线圈相互排斥，电流方向相反，B 错误；当天平示数为负时，两线圈相互吸引，电流方向相同，A 正确；线圈 Ⅰ 对线圈 Ⅱ 的作用力与线圈 Ⅱ 对线圈 Ⅰ 的作用力是一对相互作用力，等大反向，C 错误；静止时线圈 Ⅱ 平衡，线圈 Ⅰ 对线圈 Ⅱ 的作用力与托盘对线圈 Ⅱ 的作用力是一对平衡力，D 错误。
}

\item
%% number: 8
%% paperName: 2018年11月浙江选考第 8 题 3 分
%% typeId: 1
%% type: 单选题
%% chapter: 电场
%% point: 带电粒子的平衡
%% method: 无
%% score: 3
%% degree: 550
%% duplicateId: 0
%% body:
电荷量为 $4 \times 10^{-6}\UC$ 的小球绝缘固定在 A 点，质量为 $0.2\Ukg$、电荷量为 $-5 \times 10^{-6}\UC$ 的小球用绝缘细线悬挂，静止于 B 点。A、B 间距离为 $30\Ucm$，AB 连线与竖直方向夹角为 $60^\circ$。静电力常量为 $9.0 \times 10^{9}\UNmqCq$，小球可视为点电荷。下列图示正确的是\xzanswer{B}
\fourchoices[answer=B, ispicture=true, h=4.5cm]
{figs/08a.png}
{figs/08b.png}
{figs/08c.png}
{figs/08d.png}
%% answer: B
%% memo:
\memoanswer{
两球之间的库仑力为 $F = k\frac{q_A q_B}{r^{2}} = 9.0 \times 10^{9} \times \frac{4 \times 10^{-6} \times 5 \times 10^{-6}}{0.3^{2}}\UN = 2\UN$，小球 B 受到的重力大小为 $G_B = 2\UN$，且 $F$ 与竖直方向的夹角为 $60^\circ$，$F = G_B$，故小球 B 受到的库仑力、重力以及细线的拉力组成的矢量三角形为等边三角形，所以细线与竖直方向的夹角为 $60^\circ$，B 正确。
}

\item
%% number: 9
%% paperName: 2018年11月浙江选考第 9 题 3 分
%% typeId: 1
%% type: 单选题
%% chapter: 曲线运动
%% point: 圆周运动的应用
%% method: 无
%% score: 3
%% degree: 550
%% duplicateId: 0
%% body:
一质量为 $2.0 \times 10^{3}\Ukg$ 的汽车在水平公路上行驶，路面对轮胎的径向最大静摩擦力为 $1.4 \times 10^{4}\UN$，当汽车经过半径为 $80\Um$ 的弯道时，下列判断正确的是\xzanswer{D}

\onepicture[w=5.0cm]{figs/09.png}

\fourchoices[answer=D]
{汽车转弯时所受的力有重力、弹力、摩擦力和向心力}
{汽车转弯的速度为 $20\Ums$ 时所需的向心力为 $1.4 \times 10^{4}\UN$}
{汽车转弯的速度为 $20\Ums$ 时汽车会发生侧滑}
{汽车能安全转弯的向心加速度不超过 $7.0\Umsq$}
%% answer: D
%% memo:
\memoanswer{
汽车转弯时受到重力、地面的支持力以及地面的摩擦力，其中摩擦力充当向心力，A 错误；当最大静摩擦力充当向心力时对应的速度为临界速度，大于这个速度则发生侧滑，根据牛顿第二定律可得 $f = m\frac{v^{2}}{r}$，解得 $v = \sqrt{\frac{fr}{m}} = \sqrt{\frac{1.4 \times 10^{4} \times 80}{2.0 \times 10^{3}}}\Ums = \sqrt{560}\Ums = 20\sqrt{1.4}\Ums$，所以汽车转弯速度为 $20\Ums$ 时所需的向心力小于 $1.4 \times 10^{4}\UN$，汽车不会侧滑，BC 错误；汽车能安全转弯的向心加速度 $a = \frac{v^{2}}{r} = \frac{560}{80}\Umsq = 7\Umsq$，即不超过 $7.0\Umsq$，D 正确。
}

\item
%% number: 10
%% paperName: 2018年11月浙江选考第 10 题 3 分
%% typeId: 1
%% type: 单选题
%% chapter: 磁场
%% point: 洛伦兹力
%% method: 无
%% score: 3
%% degree: 450
%% duplicateId: 0
%% body:
磁流体发电的原理如图所示。将一束速度为 $v$ 的等离子体垂直于磁场方向喷入磁感应强度为 $B$ 的匀强磁场中，在相距为 $d$、宽为 $a$、长为 $b$ 的两平行金属板间便产生电压。如果把上、下板和电阻 $R$ 连接，上、下板就是一个直流电源的两极。若稳定时等离子体在两板间均匀分布，电阻率为 $\rho$。忽略边缘效应，下列判断正确的是\xzanswer{C}

\onepicture[w=7.0cm]{figs/10.png}

\fourchoices[answer=C]
{上板为正极，电流 $I = \frac{Bdvab}{Rab + \rho d}$}
{上板为负极，电流 $I = \frac{Bvad^{2}}{Rab + \rho d}$}
{下板为正极，电流 $I = \frac{Bdvab}{Rab + \rho d}$}
{下板为负极，电流 $I = \frac{Bvad^{2}}{Rab + \rho d}$}
%% answer: C
%% memo:
\memoanswer{
根据左手定则，正电荷受到的洛伦兹力方向向下，负电荷受到的洛伦兹力向上，因此下极板为电源的正极；根据平衡有 $qvB = q\frac{E}{d}$，解得稳定时电源的电动势 $E = Bdv$，则流过 $R$ 的电流为 $I = \frac{E}{R+r}$，而 $r = \rho\frac{d}{S}$，$S = ab$，则得电流大小为 $I = \frac{Bdvab}{Rab+\rho d}$，C 正确。
}

\item
%% number: 11
%% paperName: 2018年11月浙江选考第 11 题 3 分
%% typeId: 1
%% type: 单选题
%% chapter: 力物体的平衡
%% point: 摩擦力
%% method: 近似与估算
%% score: 3
%% degree: 450
%% duplicateId: 0
%% body:
小明在观察如图所示的沙子堆积时，发现沙子会自然堆积成圆锥体，且在不断堆积过程中，材料相同的沙子自然堆积成的圆锥体的最大底角都是相同的。小明测出这堆沙子的底部周长为 $31.4\Um$，利用物理知识测得沙子之间的摩擦因数为 0.5，估算出这堆沙子的体积最接近\xzanswer{A}

\onepicture[w=4.5cm]{figs/11.png}

\fourchoices[answer=A]
{$60\Umq$}
{$200\Umq$}
{$250\Umq$}
{$500\Umq$}
%% answer: A
%% memo:
\memoanswer{
沙堆底部周长为 $31.4\Um$，故圆锥体的底部圆半径为 $r = \frac{31.4}{2 \times 3.14}\Um = 5\Um$；对锥面上的一粒沙粒分析，当满足 $\mu mg\cos\theta = mg\sin\theta$（$\theta$ 为锥体的底角）时沙粒刚好静止，故 $\mu = \tan\theta = \frac{h}{r}$，解得圆锥体高 $h = 2.5\Um$，故圆锥体的体积约为 $V = \frac{1}{3}Sh = \frac{1}{3} \times 3.14 \times 5^{2} \times 2.5 \approx 60\Umc$，A 正确。

\onepicture[w=4.0cm]{figs/11_an.png}
}

\item
%% number: 12
%% paperName: 2018年11月浙江选考第 12 题 3 分
%% typeId: 1
%% type: 单选题
%% chapter: 曲线运动
%% point: 人造地球卫星
%% method: 无
%% score: 3
%% degree: 450
%% duplicateId: 0
%% body:
20 世纪人类最伟大的创举之一是开拓了太空的全新领域。现有一艘远离星球在太空中直线飞行的宇宙飞船，为了测量自身质量，启动推进器，测出飞船在短时间 $\Delta t$ 内速度的改变为 $\Delta v$，和飞船受到的推力 $F$（其它星球对它的引力可忽略）。飞船在某次航行中，当它飞近一个孤立的星球时，飞船能以速度 $v$，在离星球的较高轨道上绕星球做周期为 $T$ 的匀速圆周运动。已知星球的半径为 $R$，引力常量用 $G$ 表示。则宇宙飞船和星球的质量分别是\xzanswer{D}

\onepicture[w=4.5cm]{figs/12.png}

\fourchoices[answer=D]
{$\frac{F\Delta v}{\Delta t}$，$\frac{v^{2}R}{G}$}
{$\frac{F\Delta v}{\Delta t}$，$\frac{v^{3}T}{2\pi G}$}
{$\frac{F\Delta t}{\Delta v}$，$\frac{v^{2}R}{G}$}
{$\frac{F\Delta t}{\Delta v}$，$\frac{v^{3}T}{2\pi G}$}
%% answer: D
%% memo:
\memoanswer{
直线推进时，根据动量定理可得 $F\Delta t = m\Delta v$，解得飞船的质量为 $m = \frac{F\Delta t}{\Delta v}$；绕孤立星球运动时，根据 $G\frac{Mm}{r^{2}} = m\frac{4\pi^{2}}{T^{2}}r$，又 $G\frac{Mm}{r^{2}} = m\frac{v^{2}}{r}$，解得 $M = \frac{v^{3}T}{2\pi G}$，D 正确。
}

\item
%% number: 13
%% paperName: 2018年11月浙江选考第 13 题 3 分
%% typeId: 1
%% type: 单选题
%% chapter: 运动和力
%% point: 牛顿定律的应用
%% method: 无
%% score: 3
%% degree: 350
%% duplicateId: 0
%% body:
如图所示为某一游戏的局部简化示意图。D 为弹射装置，AB 是长为 $21\Um$ 的水平轨道，倾斜直轨道 BC 固定在竖直放置的半径为 $R = 10\Um$ 的圆形支架上，B 为圆形的最低点，轨道 AB 与 BC 平滑连接，且在同一竖直平面内。某次游戏中，无动力小车在弹射装置 D 的作用下，以 $v_0 = 10\Ums$ 的速度滑上轨道 AB，并恰好能冲到轨道 BC 的最高点。已知小车在轨道 AB 上受到的摩擦力为其重量的 0.2 倍，轨道 BC 光滑，则小车从 A 到 C 的运动时间是\xzanswer{A}

\onepicture[w=9.0cm]{figs/13.png}

\fourchoices[answer=A]
{$5\Us$}
{$4.8\Us$}
{$4.4\Us$}
{$3\Us$}
%% answer: A
%% memo:
\memoanswer{
设小车的质量为 $m$。小车在 AB 段做匀减速直线运动，加速度 $a_1 = \frac{f}{m} = \frac{0.2mg}{m} = 0.2g = 2\Umsq$；在 AB 段，根据动能定理可得 $-fx_{AB} = \frac{1}{2}mv_B^{2} - \frac{1}{2}mv_0^{2}$，解得 $v_B = 4\Ums$，故 $t_1 = \frac{10-4}{2}\Us = 3\Us$。小车在 BC 段，根据机械能守恒可得 $\frac{1}{2}mv_B^{2} = mgh_{CD}$，解得 $h_{CD} = 0.8\Um$；过圆形支架的圆心 O 点作 BC 的垂线，根据几何知识可得 $\frac{R}{x_{BC}} = \frac{\frac{1}{2}x_{BC}}{h_{CD}}$，解得 $x_{BC} = 4\Um$，$\sin\theta = \frac{h_{CD}}{x_{BC}} = \frac{1}{5}$，故小车在 BC 上运动的加速度为 $a_2 = g\sin\theta = 2\Umsq$，运动时间为 $t_2 = \frac{v_B}{a_2} = \frac{4}{2}\Us = 2\Us$，所以小车运动的总时间为 $t = t_1 + t_2 = 5\Us$，A 正确。
}

\item
%% number: 14
%% paperName: 2018年11月浙江选考第 14 题 2 分
%% typeId: 2
%% type: 多选题
%% chapter: 原子和原子核
%% point: 玻尔模型
%% method: 无
%% score: 2
%% degree: 650
%% duplicateId: 0
%% body:
处于较高能级的氢原子向较低能级跃迁时，能辐射出 a、b 两种可见光，a 光照射某金属表面时有光电子逸出，b 光照射该金属表面时没有光电子逸出，则\xzanswer{BD}
\fourchoices[answer=BD]
{以相同的入射角射向一平行玻璃砖，a 光的侧移量小于 b 光的}
{垂直入射到同一单缝衍射装置，a 光的衍射中央亮条纹宽度小于 b 光的}
{a 光和 b 光的频率之比可能是 $\frac{20}{27}$}
{a 光子的动量大于 b 光子的}
%% answer: BD
%% memo:
\memoanswer{
发生光电效应的条件是入射光的频率大于极限频率。a 光照射某金属表面时有光电子逸出，b 光照射时没有光电子逸出，故 a 光的频率大于 b 光的频率，C 错误；因为 a 光的频率大于 b 光的频率，在同种介质中 a 光的折射率大、偏折程度大、侧移量大，A 错误；频率越大波长越短，故 a 光的波长较短，由光的波长越长衍射中央亮条纹越宽知，垂直入射到同一单缝衍射装置时 a 光的衍射中央亮条纹宽度小于 b 光的，B 正确；由 $p = \frac{h}{\lambda}$ 知，a 光子的动量大于 b 光子的，D 正确。
}

\item
%% number: 15
%% paperName: 2018年11月浙江选考第 15 题 2 分
%% typeId: 2
%% type: 多选题
%% chapter: 原子和原子核
%% point: 人工核反应
%% method: 无
%% score: 2
%% degree: 550
%% duplicateId: 0
%% body:
一个铍原子核（$\ce{^7_4 Be}$）俘获一个核外电子（通常是最靠近原子核的 K 壳层的电子）后发生衰变，生成一个锂核（$\ce{^7_3 Li}$），并放出一个不带电的质量接近零的中微子 $\nu_e$，人们把这种衰变称为“K 俘获”。静止的铍核发生了“K 俘获”，其核反应方程为

\ce{^7_4 Be + ^0_{-1} e -> ^7_3 Li + \nu_e}

已知铍原子的质量为 $M_{Be} = 7.016929\Uu$，锂原子的质量为 $M_{Li} = 7.016004\Uu$，1 u 相当于 $9.31 \times 10^{2}\UMeV$。下列说法正确的是\xzanswer{AC}
\fourchoices[answer=AC]
{中微子的质量数和电荷数均为零}
{锂核（$\ce{^7_3 Li}$）获得的动能约为 $0.86\UMeV$}
{中微子与锂核（$\ce{^7_3 Li}$）的动量之和等于反应前电子的动量}
{中微子与锂核（$\ce{^7_3 Li}$）的能量之和等于反应前电子的能量}
%% answer: AC
%% memo:
\memoanswer{
中微子不带电且质量接近零，故其质量数和电荷数均为零，A 正确；锂核获得的动能约为核衰变时质量亏损所释放出的能量 $E = 9.31 \times 10^{2}(M_{Be} + M_{e} - M_{Li})\UMeV > 0.86\UMeV$，故 B 错误；反应过程中系统内力远大于外力，系统动量守恒，由于反应前铍核静止，故中微子与锂核的动量之和等于反应前电子的动量，C 正确；由能量守恒知中微子与锂核的能量之和不等于反应前电子的能量，D 错误。
}

\item
%% number: 16
%% paperName: 2018年11月浙江选考第 16 题 2 分
%% typeId: 2
%% type: 多选题
%% chapter: 机械波
%% point: 波的图象
%% method: 无
%% score: 2
%% degree: 550
%% duplicateId: 0
%% body:
如图所示，两种不同材料的弹性细绳在 O 处连接，M、O 和 N 是该绳上的三个点，OM 间距离为 $7.0\Um$，ON 间距离为 $5.0\Um$。O 点上下振动，则形成以 O 点为波源向左和向右传播的简谐横波 Ⅰ 和 Ⅱ，其中波 Ⅱ 的波速为 $1.0\Ums$。$t=0$ 时刻 O 点处波谷位置，观察发现 $5\Us$ 后此波谷传到 M 点，此时 O 点正通过平衡位置向上运动，OM 间还有一个波谷。则\xzanswer{BD}

\onepicture[w=8.0cm]{figs/16.png}

\fourchoices[answer=BD]
{波 Ⅰ 的波长为 $4\Um$}
{N 点的振动周期为 $4\Us$}
{$t=3\Us$ 时，N 点恰好处于波谷}
{当 M 点处于波峰时，N 点也一定处于波峰}
%% answer: BD
%% memo:
\memoanswer{
由“$5\Us$ 后此波谷传到 M 点，此时 O 点正通过平衡位置向上运动”可知，OM 之间的距离为 $1\frac{1}{4}\lambda_1 = 7\Um$，解得波 Ⅰ 的波长 $\lambda_1 = 5.6\Um$，故 A 错误；波 Ⅰ 的波速 $v_1 = \frac{7.0}{5}\Ums = 1.4\Ums$，周期 $T = \frac{\lambda_1}{v_1} = 4\Us$，同一波源的频率相同，故 N 点的振动周期为 $4\Us$，B 正确；波 Ⅱ 的波长为 $\lambda_2 = v_2T = 4\Um$，$t=0$ 时刻 N 点处于平衡位置向下振动，经过 $3\Us$（即四分之三周期），N 点在波峰，C 错误；因为 M、N 两点到波源的距离都为各自波长的 $1\frac{1}{4}$，且振动周期相同、起振方向相同，所以两者振动步调相同，当 M 点处于波峰时 N 点也一定处于波峰，D 正确。
}

\item
%% number: 17
%% paperName: 2018年11月浙江选考第 17 题 2 分
%% typeId: 4
%% type: 实验题
%% chapter: 运动和力
%% point: 验证牛顿第二定律
%% method: 无
%% score: 2
%% degree: 650
%% duplicateId: 0
%% body:
\begin{enumerate}
	\item
	在“探究求合力的方法”的实验中，下列操作正确的是（多选）

	\fourchoices[answer=AD]
	{在使用弹簧秤时，使弹簧秤与木板平面平行}
	{每次拉伸橡皮筋时，只要使橡皮筋伸长量相同即可}
	{橡皮筋应与两绳夹角的平分线在同一直线上}
	{描点确定拉力方向时，两点之间的距离应尽可能大一些}
	\item
	在“探究加速度与力、质量的关系”的实验中，两个相同的小车放在光滑水平板上，前端各系一条细绳，绳的另一端跨过定滑轮各挂一个小盘，盘中可放重物。小车的停和动通过用黑板擦按住小车后的细线和抬起来控制，如图 \ref{201811月浙江17a} 所示。实验要求小盘和重物所受的重力近似等于使小车做匀加速直线运动的力。
	\begin{steps}
		\item
		请指出图 \ref{201811月浙江17b} 中错误之处：\tkanswer[6.0cm]{拉小车的细绳与木板没有平行，托盘和砝码的总质量 $m$ 没有远小于小车的质量 $M$}。
		\item
		调整好装置后，在某次实验中测得两小车的位移分别是 $x_{1}$ 和 $x_{2}$，则两车的加速度之比为 \tkanswer[2.0cm]{$x_{1}:x_{2}$}。
	\end{steps}
\end{enumerate}

\twopicture[numA=1, numB=2, numstyle=arabic, labelA=201811月浙江17a, labelB=201811月浙江17b, align=right, wA=6.0cm, wB=6.0cm]
{figs/17a.png}
{figs/17b.png}

%% answer:
\jdanswer{
\begin{enumerate}
	\item
	AD

	\item
	① 拉小车的细绳与木板没有平行，托盘和砝码的总质量 $m$ 没有远小于小车的质量 $M$；② $x_{1}:x_{2}$

\end{enumerate}
}
%% memo:
\memoanswer{
（1）为减小实验误差，应使弹簧秤与木板平面平行，A 正确；为了使力的作用效果相同，每次拉伸橡皮筋时应使橡皮筋的伸长方向和伸长量都相同，即结点在同一位置，而不是只要伸长量相同，也不一定使橡皮筋在两绳夹角的平分线上，BC 错误；描点确定拉力方向时，为了减小实验误差，两点间的距离应适当大些，D 正确。

（2）① 细绳拉力方向应与木板平行，以保证小车做匀加速直线运动的合力等于绳子的拉力；为使小车受到的拉力约等于托盘和砝码的重力，托盘和砝码的总质量应远小于小车的质量。故错误之处为：拉小车的细绳与木板没有平行，托盘和砝码的总质量 $m$ 没有远小于小车的质量 $M$。

② 两车同时释放、同时静止，运动时间相同，根据 $x = \frac{1}{2}at^{2}$ 知，加速度之比等于位移之比，即 $\frac{a_1}{a_2} = \frac{x_1}{x_2}$。
}

\item
%% number: 18
%% paperName: 2018年11月浙江选考第 18 题 5 分
%% typeId: 4
%% type: 实验题
%% chapter: 电路
%% point: 伏安法测电阻
%% method: 无
%% score: 5
%% degree: 550
%% duplicateId: 0
%% body:
为了比较精确地测定阻值未知的定值电阻 $R_{x}$，小明设计了如图 \ref{201811月浙江18a} 所示的电路。
\begin{enumerate}
	\item
	实验时，闭合开关 S，滑动变阻器的滑片滑至合适位置保持不变，将 c 点先后与 a、b 点连接，发现电压表示数变化较大，电流表示数基本不变，则测量时应将 c 点接 \tkanswer[1.2cm]{a 点}（选填“a 点”或“b 点”），按此连接测量，测量结果 \tkanswer[1.2cm]{小于}（选填“小于”、“等于”或“大于”）$R_{x}$ 的真实值。
	\item
	根据实验测得的 6 组数据，在图 \ref{201811月浙江18b} 中描点，作出了 2 条图线。你认为正确的是 \tkanswer[0.8cm]{②}（选填“①”或“②”），并由图线求出电阻 $R_{x} =$ \tkanswer[1.2cm]{7.5}$\UO$。（保留两位有效数字）
\end{enumerate}

\twopicture[numA=1, numB=2, numstyle=arabic, labelA=201811月浙江18a, labelB=201811月浙江18b, align=right, wA=4.0cm, wB=8.0cm]
{figs/18a.png}
{figs/18b.png}

%% answer:
\jdanswer{
\begin{enumerate}
	\item
	a 点，小于

	\item
	②，7.5

\end{enumerate}
}
%% memo:
\memoanswer{
（1）c 先后与 a、b 连接，电压表示数变化较大、电流表示数基本不变，说明电流表的内阻与被测电阻的阻值相差不大，故电流表应采用外接法，即 c 接 a 点；由于电压表的分流，电流的测量值大于通过被测电阻的真实值，而电压的测量值等于被测电阻的真实值，由欧姆定律知测量值小于真实值。

（2）因为电压表测量的是被测电阻两端的电压，当其两端电压为零时电流也为零，图线应过原点，故②正确；图线的斜率表示电阻的倒数，由题图知 $R_{x} = \frac{3.00}{0.40}\UO = 7.5\UO$。
}

\item
%% number: 19
%% paperName: 2018年11月浙江选考第 19 题 9 分
%% typeId: 6
%% type: 计算题
%% chapter: 曲线运动
%% point: 平抛运动
%% method: 无
%% score: 9
%% degree: 650
%% duplicateId: 0
%% body:
在竖直平面内，某一游戏轨道由直轨道 AB 和弯曲的细管道 BCD 平滑连接组成，如图所示。小滑块以某一初速度从 A 点滑上倾角为 $\theta = 37^\circ$ 的直轨道 AB，到达 B 点的速度大小为 $2\Ums$，然后进入细管道 BCD，从细管道出口 D 点水平飞出，落到水平面上的 G 点。已知 B 点的高度 $h_{1} = 1.2\Um$，D 点的高度 $h_{2} = 0.8\Um$，D 点与 G 点间的水平距离 $L = 0.4\Um$，滑块与轨道 AB 间的动摩擦因数 $\mu = 0.25$，$\sin37^\circ = 0.6$，$\cos37^\circ = 0.8$。
\begin{enumerate}
	\item
	求小滑块在轨道 AB 上的加速度和在 A 点的初速度；
	\item
	求小滑块从 D 点飞出的速度；
	\item
	判断细管道 BCD 的内壁是否光滑。
\end{enumerate}

\onepicture[w=8.0cm, align=right]{figs/19.png}

%% answer:
\jdanswer{
\begin{enumerate}
	\item
	$8\Umsq$，$6\Ums$

	\item
	$1\Ums$

	\item
	小滑块动能减小，重力势能也减小，机械能不守恒，所以细管道 BCD 内壁不光滑。

\end{enumerate}
}
%% memo:
\memoanswer{
（1）小滑块在上滑过程中，由牛顿第二定律得
$mg\sin\theta + \mu mg\cos\theta = ma$，
解得 $a = 8\Umsq$，方向沿斜面向下；由运动学公式得
$v_B^{2} - v_0^{2} = -2a\frac{h_1}{\sin\theta}$，解得 $v_0 = 6\Ums$。

（2）小滑块在 D 处水平飞出，做平抛运动，由
$L = v_D t$，$h_2 = \frac{1}{2}gt^{2}$，
解得 $v_D = 1\Ums$。

（3）小滑块动能减小，重力势能也减小，机械能不守恒，所以细管道 BCD 内壁不光滑。
}

\item
%% number: 20
%% paperName: 2018年11月浙江选考第 20 题 12 分
%% typeId: 6
%% type: 计算题
%% chapter: 运动和力
%% point: 变加速直线运动
%% method: 无
%% score: 12
%% degree: 550
%% duplicateId: 0
%% body:
如图所示，在地面上竖直固定了刻度尺和轻质弹簧，弹簧原长时上端与刻度尺上的 A 点等高。质量 $m = 0.5\Ukg$ 的篮球静止在弹簧正上方，其底端距 A 点高度 $h_1 = 1.10\Um$。篮球静止释放，测得第一次撞击弹簧时，弹簧的最大形变量 $x_1 = 0.15\Um$，第一次反弹至最高点，篮球底端距 A 点的高度 $h_2 = 0.873\Um$，篮球多次反弹后静止在弹簧的上端，此时弹簧的形变量 $x_2 = 0.01\Um$，弹性势能为 $E_p = 0.025\UJ$。若篮球运动时受到的空气阻力大小恒定，忽略篮球与弹簧碰撞时的能量损失和篮球的形变，弹簧形变在弹性限度范围内。求：
\begin{enumerate}
	\item
	弹簧的劲度系数；
	\item
	篮球在运动过程中受到的空气阻力；
	\item
	篮球在整个运动过程中通过的路程；
	\item
	篮球在整个运动过程中速度最大的位置。
\end{enumerate}

\onepicture[w=5.0cm, align=right]{figs/20.png}

%% answer:
\jdanswer{
\begin{enumerate}
	\item
	$500\UNm$

	\item
	$0.5\UN$

	\item
	$11.05\Um$

	\item
	$0.009\Um$

\end{enumerate}
}
%% memo:
\memoanswer{
（1）篮球静止在弹簧上，由平衡条件得
$mg - kx_{2} = 0$，解得 $k = 500\UNm$。

（2）篮球从开始运动到第一次上升到最高点，由动能定理得
$mg(h_{1}-h_{2}) - f(h_{1}+h_{2}+2x_{1}) = 0$，
解得 $f = 0.5\UN$。

（3）设球在整个运动过程中通过的总路程为 $s$，由能量守恒得
$mg(h_{1}+x_{2}) = fs + E_{p}$，解得 $s = 11.05\Um$。

（4）球在首次下落过程中，合力为零处速度最大，设速度最大时弹簧形变量为 $x_{3}$，由平衡条件得
$mg - f - kx_{3} = 0$，解得 $x_{3} = 0.009\Um$，
即在 A 点下方、离 A 点 $0.009\Um$ 处。
}

\item
%% number: 21
%% paperName: 2018年11月浙江选考第 21 题 4 分
%% typeId: 4
%% type: 实验题
%% chapter: 运动和力
%% point: 动量守恒定律
%% method: 无
%% score: 4
%% degree: 550
%% duplicateId: 0
%% body:
小明做“探究碰撞中的不变量”实验的装置如图 \ref{201811月浙江21a} 所示，悬挂在 O 点的单摆由长为 $l$ 的细线和直径为 $d$ 的小球 A 组成，小球 A 与放置在光滑支撑杆上的直径相同的小球 B 发生对心碰撞，碰后小球 A 继续摆动，小球 B 做平抛运动。
\begin{enumerate}
	\item
	小明用游标卡尺测小球 A 直径如图 \ref{201811月浙江21b} 所示，则 $d =$ \tkanswer[1.4cm]{14.40}$\Umm$。又测得了小球 A 质量 $m_{1}$、细线长度 $l$、碰撞前小球 A 拉起的角度 $\alpha$ 和碰撞后小球 B 做平抛运动的水平位移 $x$、竖直下落高度 $h$。为完成实验，还需要测量的物理量有：\tkanswer[4.5cm]{小球 B 质量 $m_{2}$、碰后小球 A 摆动的最大角 $\beta$}。
	\item
	若 A、B 两球碰后粘在一起形成新单摆，其周期 \tkanswer[1.0cm]{大于}（选填“小于”、“等于”或“大于”）粘合前单摆的周期（摆角小于 $5^\circ$）。
\end{enumerate}

\twopicture[numA=1, numB=2, numstyle=arabic, labelA=201811月浙江21a, labelB=201811月浙江21b, align=right, wA=4.0cm, wB=7.0cm]
{figs/21a.png}
{figs/21b.png}

%% answer:
\jdanswer{
\begin{enumerate}
	\item
	14.40，小球 B 质量 $m_{2}$、碰后小球 A 摆动的最大角 $\beta$

	\item
	大于

\end{enumerate}
}
%% memo:
\memoanswer{
（1）由题图知，游标卡尺的精确度为 $0.05\Umm$，则小球的直径 $d = 14\Umm + 8 \times 0.05\Umm = 14.40\Umm$。A 球从释放到刚运动到最低点的过程中，由动能定理得 $m_1gl(1 - \cos\alpha) = \frac{1}{2}m_1v_1^{2}$，可得出 $v_1$；由平抛运动规律得 $h = \frac{1}{2}gt^{2}$、$x = v_2t$，可得出 $v_2$。A、B 碰撞过程中动量守恒，根据动量守恒定律得 $m_1v_1 = m_1v_1' + m_2v_2'$，因此还需要测量 $m_2$ 和碰后 A 球的速度，其中碰后 A 球的速度通过 A 球向右摆动的最大摆角来计算，故还需要测量的物理量是 $m_2$ 和 A 球向右的最大摆角。

（2）若 A、B 两球碰后粘在一起形成新单摆，此时重心位置变为两球中间位置，悬点到新重心的距离变长，即摆长变长，根据单摆周期公式 $T = 2\pi\sqrt{\frac{l}{g}}$ 知，单摆的周期变大。
}

\item
%% number: 22
%% paperName: 2018年11月浙江选考第 22 题 10 分
%% typeId: 6
%% type: 计算题
%% chapter: 电磁感应
%% point: 电磁感应含容电路
%% method: 无
%% score: 10
%% degree: 450
%% duplicateId: 0
%% body:
如图所示，在间距 $L = 0.2\Um$ 的两光滑平行水平金属导轨间存在方向垂直于纸面（向内为正）的磁场，磁感应强度分布沿 $y$ 方向不变，沿 $x$ 方向如下：

\[B = \left\{ \begin{array}{l} 1\mathrm{~T}, \quad x > 0.2\mathrm{~m} \\ 5x\mathrm{~T}, \quad -0.2\mathrm{~m} \le x \le 0.2\mathrm{~m} \\ -1\mathrm{~T}, \quad x < -0.2\mathrm{~m} \end{array} \right.\]

导轨间通过单刀双掷开关 S 连接恒流源和电容 $C = 1\UF$ 的未充电的电容器，恒流源可为电路提供恒定电流 $I = 2\UA$，电流方向如图所示。有一质量 $m = 0.1\Ukg$ 的金属棒 ab 垂直导轨静止放置于 $x_0 = 0.7\Um$ 处。开关 S 掷向 1，棒 ab 从静止开始运动，到达 $x_3 = -0.2\Um$ 处时，开关 S 掷向 2。已知棒 ab 在运动过程中始终与导轨垂直。求：

（提示：可以用 $F-x$ 图象下的“面积”代表力 $F$ 所做的功）
\begin{enumerate}
	\item
	棒 ab 运动到 $x_{1} = 0.2\Um$ 时的速度 $v_{1}$；
	\item
	棒 ab 运动到 $x_{2} = -0.1\Um$ 时的速度 $v_{2}$；
	\item
	电容器最终所带的电荷量 $Q$。
\end{enumerate}

\onepicture[w=8.0cm, align=right]{figs/22.png}

%% answer:
\jdanswer{
\begin{enumerate}
	\item
	$2\Ums$

	\item
	$\sqrt{4.6}\Ums$

	\item
	$\frac{2}{7}\UC$

\end{enumerate}
}
%% memo:
\memoanswer{
（1）金属棒受到的安培力为
$F = BIL$，
由牛顿第二定律得
$a = \frac{BIL}{m}$，
由运动学公式知金属棒的速度为
$v_{1} = \sqrt{2a(x_{0} - x_{1})} = 2\Ums$。

（2）在区间 $-0.2\Um \le x \le 0.2\Um$ 内，金属棒受到的安培力为
$F = 5xIL$，
作出 $F-x$ 图象，如图所示。由图象知安培力做的功为
$W = \frac{5IL}{2}(x_{1}^{2} - x_{2}^{2})$，
由动能定理得
$W = \frac{1}{2}mv_{2}^{2} - \frac{1}{2}mv_{1}^{2}$，解得 $v_{2} = \sqrt{4.6}\Ums = \frac{\sqrt{460}}{10}\Ums$。

（3）金属棒运动的过程中，由动量定理得
$-BLQ = mv - mv_{3}$，
又电荷量
$Q = CBLv$，
又因 $x = -0.2\Um$ 处的速度
$v_{3} = v_{1} = 2\Ums$，
联立解得 $Q = \frac{CBLmv_{3}}{CB^{2}L^{2} + m} = \frac{2}{7}\UC$。

\onepicture[w=5.0cm]{figs/22_an.png}
}

\item
%% number: 23
%% paperName: 2018年11月浙江选考第 23 题 10 分
%% typeId: 6
%% type: 计算题
%% chapter: 磁场
%% point: 带电粒子在磁场中的运动
%% method: 分情况讨论
%% score: 10
%% degree: 250
%% duplicateId: 0
%% body:
小明受回旋加速器的启发，设计了如图 \ref{201811月浙江23a} 所示的“回旋变速装置”。两相距为 $d$ 的平行金属栅极板 M、N，板 M 位于 $x$ 轴上，板 N 在它的正下方。两板间加上如图 \ref{201811月浙江23b} 所示的幅值为 $U_{0}$ 的交变电压，周期 $T_{0} = \frac{2\pi m}{qB}$。板 M 上方和板 N 下方有磁感应强度大小均为 $B$、方向相反的匀强磁场。粒子探测器位于 $y$ 轴处，仅能探测到垂直射入的带电粒子。

有一沿 $x$ 轴可移动、粒子出射初动能可调节的粒子发射源，沿 $y$ 轴正方向射出质量为 $m$、电荷量为 $q$（$q > 0$）的粒子。$t = 0$ 时刻，发射源在 $(x, 0)$ 位置发射一带电粒子。忽略粒子的重力和其它阻力，粒子在电场中运动的时间不计。
\begin{enumerate}
	\item
	若粒子只经磁场偏转并在 $y = y_{0}$ 处被探测到，求发射源的位置和粒子的初动能；
	\item
	若粒子两次进出电场区域后被探测到，求粒子发射源的位置 $x$ 与被探测到的位置 $y$ 之间的关系。
\end{enumerate}

\twopicture[numA=1, numB=2, numstyle=arabic, labelA=201811月浙江23a, labelB=201811月浙江23b, align=right, wA=6.0cm, wB=6.0cm]
{figs/23a.png}
{figs/23b.png}

%% answer:
\jdanswer{
\begin{enumerate}
	\item
	$x_{0} = y_{0}$，$E_{k0} = \frac{(qBy_{0})^{2}}{2m}$

	\item
	（i）若 $E_{k0} > 2qU_{0}$，$x = y + \frac{2}{qB}\sqrt{(yqB)^{2} + 2mqU_{0}} + \frac{2}{qB}\sqrt{(yqB)^{2} + 4mqU_{0}}$；

	（ii）若 $qU_{0} < E_{k0} < 2qU_{0}$，$x = -3(y + d) + \frac{2}{qB}\sqrt{(y + d)^{2}q^{2}B^{2} + 2mqU_{0}}$；

	（iii）若 $E_{k0} < qU_{0}$，$x = -y - d + \frac{4}{qB}\sqrt{(y + d)^{2}q^{2}B^{2} - 2mqU_{0}}$。

\end{enumerate}
}
%% memo:
\memoanswer{
（1）由题意知，粒子运动轨迹为四分之一圆弧，由几何关系知轨迹的半径为
$R = y_{0}$，
因此发射源的位置 $x_{0} = y_{0}$；由洛伦兹力提供向心力得
$qvB = m\frac{v^{2}}{R}$，
又 $E_{k0} = \frac{1}{2}mv^{2}$，
故粒子的初动能 $E_{k0} = \frac{(qBy_{0})^{2}}{2m}$。

（2）分下面三种情况讨论：

① 如图 1，当 $E_{k0} > 2qU_{0}$ 时，有
$y = \frac{mv_{2}}{qB}$，$R_{0} = \frac{mv_{0}}{qB}$，$R_{1} = \frac{mv_{1}}{qB}$，
由动能定理得
$\frac{1}{2}mv_{1}^{2} = \frac{1}{2}mv_{0}^{2} - qU_{0}$，
$\frac{1}{2}mv_{2}^{2} = mv_{1}^{2} - qU_{0}$，
发射源的位置为
$x = y + 2(R_{0} + R_{1})$，
解得
$x = y + \frac{2}{qB}\sqrt{(yqB)^{2} + 2mqU_{0}} + \frac{2}{qB}\sqrt{(yqB)^{2} + 4mqU_{0}}$。

\onepicture[w=5.0cm]{figs/23_an1.png}

② 如图 2，当 $qU_{0} < E_{k0} < 2qU_{0}$ 时，有
$-y - d = \frac{mv_{2}}{qB}$，$R_{0} = \frac{mv_{0}}{qB}$，
由动能定理得
$\frac{1}{2}mv_{0}^{2} = \frac{1}{2}mv_{2}^{2} + qU_{0}$，
发射源的位置为
$x = 3(-y - d) + 2R_{0}$，
解得 $x = -3(y + d) + \frac{2}{qB}\sqrt{(y + d)^{2}q^{2}B^{2} + 2mqU_{0}}$。

\onepicture[w=5.0cm]{figs/23_an2.png}

③ 如图 3，当 $E_{k0} < qU_{0}$ 时，有
$-y - d = \frac{mv_{2}}{qB}$，$R_{0} = \frac{mv_{0}}{qB}$，
由动能定理得
$\frac{1}{2}mv_{0}^{2} = \frac{1}{2}mv_{2}^{2} - qU_{0}$，
发射源的位置为
$x = -y - d + 4R_{0}$，
解得
$x = -y - d + \frac{4}{qB}\sqrt{(y + d)^{2}q^{2}B^{2} - 2mqU_{0}}$。

\onepicture[w=5.0cm]{figs/23_an3.png}
}

\end{enumerate}

\end{document}
